\chapter{观测对象与问题的确定}
\section{先确定观测量，再选择算法}
质谱仪输出的是检测器的响应，不是未经变换的物质生成速率。对同一组数据，可以提出质量峰恢复、混合谱分离、时间响应反演和绝对定量四种不同问题。它们的未知量、先验与验收条件不同，不能以同一个“反卷积”名称掩盖这些差别。

\begin{definition}[质谱观测矩阵]
取质量通道 $m_1,\ldots,m_M$ 和时间样点 $t_1,\ldots,t_N$，定义
\begin{equation}
Y_{mn}=y_{m_m}(t_n),\qquad Y\in\R^{M\times N}.
\end{equation}
行代表质量通道，列代表采样时间。每一行是一条离子流曲线，每一列是一个时刻的采样谱。
\end{definition}
该矩阵可以画为热图，也可逐行画曲线。图~\ref{fig:01_mass_time} 是本报告使用的合成观测，其二维性来自两种实验坐标，而不是来自图像的两个空间坐标。质谱数据还可能具有保留时间、离子迁移时间等维度；多维数据组织并不意味着应采用空间图像模型\cite{deimos}。
\fig{01_mass_time}{.93}{合成的质量通道--时间观测矩阵。质量通道为离散类别，图中相邻两行不必对应相邻整数质量，更不表示两个空间像素。}

\section{三种线性变换不能混为一谈}
质量轴线形使一个窄质量峰展宽，数学上可在局部近似为沿质量轴的卷积。碎裂和不同分子的重叠谱构成通道混合，通常是矩阵乘法。膜、管路和采样腔的传输过程则使生成速率延迟并展宽，构成时间卷积。三者的结构为
\begin{equation}
\text{源项}\ \xrightarrow{\text{时间传输}}\ \text{组分响应}
\ \xrightarrow{\text{谱混合}}\ \text{质量通道信号},
\end{equation}
而质量轴展宽是检测器谱分辨率所对应的另一层变换。
本软件因此提供 \code{spectral}、\code{spectrotemporal} 和独立的质量轴 \code{mass-profile} 路径。第一条路径的估计量是经过传输后的组分响应，只有第二条路径在已知传输核条件下才尝试恢复源项。

\section{数据合同}
一个数值矩阵只有在坐标、单位、有效性掩码与实验条件明确后才成为可解释的观测。\code{MSDataset} 强制保存以下信息：时间数组 \code{time_s}、有序质量通道 \code{mz}、以通道为行的 \code{signal}、\code{signal_unit} 和 \code{time_origin}。电化学电流使用自己的 \code{ec_time_s} 与 \code{current_A}，不能为了方便而先插值到质谱时间轴。

同一时间轴必须严格递增；重复时间、无穷大、NaN、维度不匹配与重复质量通道均拒绝导入。非均匀时间不是错误数据，但它不能直接进入本版均匀 FFT 的时间反卷积路径。每个通道可以提供有效性掩码；掩码位置不参加拟合，也不作为零值观测。至少应有一个有效样点。

负离子流或负背景残差不在导入时删除。它可能来自基线定义、噪声、仪器偏移或真实的信号减小。非负约束只能施加于有相应物理意义的未知源项，不能据此把原始负数改为零。

\begin{example}[单位错误对正则化的影响]
同一电流用 A 与 nA 表示，数值相差 $10^9$。未加权残差平方相差 $10^{18}$；保持同一个数值的 L1 系数并不能表示同一个先验。若同时把信号和标准差转换为相同单位，白化残差才保持不变。即使如此，源项先验及核的单位仍需相应转换。
\end{example}

\section{时间同步与顺序扫描}
绝对时刻、相对起点和仪器内部时钟偏移是三个不同概念。\code{time_origin} 是数据生产者对时间坐标的明确声明，软件不会把相关峰最高的位置自动当作零时差。采样传输滞后属于物理模型，也不能一律当作时钟误差。

顺序扫描的各质量通道可能在一帧内不同时间采样。严格模型应写为
\begin{equation}
Y_{mn}=\sum_jS_{mj}(h_j*x_j)(t_n+\delta_m)+b_m(t_n+\delta_m)+\varepsilon_{mn}.
\end{equation}
当前 MCP 路径仅接受共同时间网格；非零 \code{channel_offsets_s} 会被明确拒绝，而不是忽略。一般采样算子可在 Python 层组合实现，但本版没有把任意时钟校正包装为已验证的自动功能。

\section{计算结论的层级}
本报告区分三层结论。第一层是数值结论：算法是否到达停止条件，前向重构与观测是否相符。第二层是模型结论：给定候选组分、背景和传输核时，解是否可辨识，残差是否包含结构。第三层是实验结论：候选分子是否得到独立证据支持，以及标定能否把响应转换为绝对物质量。

前一层成立不能推出后一层成立。一个错误的候选谱也可能拟合有限通道；一个条件良好的响应矩阵也不能证明样品中不存在未列入的组分。程序因此保存输入标签但不把标签当作化学鉴定结果；在结果中分别返回数值状态、科学解释状态与定量状态。
